% ===================================================================
%  TABELUL Nr. 7 — Satul iarna. — Ferma primăvara.
%  Traduction 2026 d'apres texto/io, controlee sur texto/fr, texto/en
%  et texto/es. Consigne : voir texto/ro/10-tabelul-01.tex.
%
%  OUVERTURE DE LA DEUXIEME SERIE : « A DOUA SERIE ».
%
%  LE TABLEAU DE LA FERME EST CELUI OU LE ROUMAIN EST LE PLUS LATIN DE
%  TOUT LE VOLUME, et c'est le contraire de ce qu'on attendrait. Le
%  vocabulaire pastoral est reste romain presque entier : « oaie » la
%  brebis, « berbec » le belier, « miel » l'agneau, « capra » la
%  chevre, « vaca » la vache, « bou » le bœuf, « cal » le cheval,
%  « pastor » le berger, « lapte » le lait, « caș » le fromage. Dix
%  mots du meme troupeau, dix mots latins. C'est la vie pastorale qui a
%  porte la langue a travers les siecles ou elle n'avait pas d'ecrit.
% ===================================================================

%%K t07-noto-1 noto
\VUnotes{143.2mm}{%
(*) Amănuntele mesei vezi în tabelele 6 și 13.%
}

%%K t07-apar-1 apar
\VUsaut{4.0mm}
\VUcentre{13.2pt}{90}{A DOUA SERIE}
\VUsaut{3.0mm}
\VUfilet{16mm}
\VUsaut{5.6mm}
\VUtitre{15.0pt}{60}{TABELUL Nr. 7}
\VUsaut{3.0mm}

%%K t07-tit-1 sub
\VUcentre{12.0pt}{0}{\textit{Prima scenă.}}
\VUsaut{3.0mm}

%%K t07-tit-2 sub
\VUpk{10.2pt}{40}{Satul iarna.}
\VUsaut{4.0mm}

%%K t07-c1-01-1 p
1. --- De sărbătorile anului nou am mers cu tata la Satul Nou, un cătun unde el
avea treburi.

%%K t07-c1-02-1 p
2. --- Am mers cu \VUgras{diligența} \textsuperscript{(11)}, care umblă între oraș
și acest sat; dar fiindcă era foarte frig, călătoria într-o trăsură atât de
neîncăpătoare nu a fost deloc plăcută.

%%K t07-c1-03-1 p
3. --- De aceea ne-am bucurat din inimă când \VUgras{surugiul} și-a oprit diligența
în piață, înaintea \VUgras{hanului} \textsuperscript{(2)} \textit{«Ursul Alb»}.
\VUgras{Hangiul} \textsuperscript{(4)} ne aștepta în prag, fumându-și liniștit
\VUgras{pipa}.

%%K t07-c1-03-2 p
Ne-a poftit înăuntru, și nu a trebuit să fiu rugat de două ori ca să mă așez
înaintea focului de viță care trosnea în vatră. Atunci nu mi-a mai păsat de
\VUgras{crivățul} tăios, de care scârțâia vechea \VUgras{firmă}
\textsuperscript{(3)} și care ducea în rotocoale negre \VUgras{fumul}
\textsuperscript{(1)} ce ieșea din coș.

%%K t07-c1-04-1 p
4. --- Era ceasul mesei. Fără zăbavă am dat cinstea cuvenită mâncării simple, dar
gustoase, care ni s-a dat (*).

%%K t07-tit-3 sub
\VUpk{10.2pt}{40}{Câteva vizite.}
\VUsaut{3.0mm}

%%K t07-c1-05-1 p
5. --- După masă am mers cu tata, care este agent de asigurări, la clienții lui.
Mai întâi l-am cercetat pe \VUgras{fierar} \textsuperscript{(5)}, care locuiește
lângă han. Potcovea un cal. M-am minunat de puterea ucenicului său, care ținea
strâns \VUgras{copita} calului lângă șorțul lui de piele, în vreme ce fierarul
bătea \VUgras{potcoava} \textsuperscript{(6)} cu lovituri grele de ciocan.

%%K t07-c1-06-1 p
6. --- Apoi am mers la \VUgras{cizmarul} satului \textsuperscript{(7)}, la
bătrânul Iacob. Deși îi slujea drept firmă o \VUgras{cizmă} minunată, se
mulțumea să petice o pereche veche de ghete scâlciate.

%%K t07-c1-06-2 p
Bătrânul ne-a spus râzând: «La oraș eram „cizmar”, și nu aveam mușterii. Aici sunt
„cârpaci”, și am de lucru cât poftești».

%%K t07-c1-07-1 p
7. --- Ne-a vorbit despre vecinul său \VUgras{bărbierul} \textsuperscript{(8)}, ai
cărui mușterii sunt toți nemulțumiți. \VUgras{Briciurile} lui sunt totdeauna
știrbite, iar \VUgras{foarfecele} nu taie niciodată cum trebuie. Dar fiindcă este
singurul bărbier din sat, oamenilor, se înțelege, nu le rămâne decât să se
bărbierească și să se tundă la el. Pe deasupra este un om zgârcit și neprietenos.

%%K t07-c1-08-1 p
8. --- Din fericire, ceilalți \VUgras{vecini} nu îi seamănă. \VUgras{Rotarul}
\textsuperscript{(9)}, de pildă, este un om bun, harnic și îndatoritor. Să îi dai
căruța la dres este o plăcere. Lucrul lui este totdeauna curat făcut.

%%K t07-c1-09-1 p
9. --- Nici de \VUgras{șelar} \textsuperscript{(10)} nu s-a plâns bătrânul Iacob,
căruia uneori îi ajută să coasă \VUgras{hamurile}, când munca nu rabdă zăbavă.

%%K t07-c1-09-2 p
Tata l-a întrebat de familie. «Toți sănătoși. Nepoata mea este la \VUgras{școală}
\textsuperscript{(12)}. \VUgras{Învățătoarea} \textsuperscript{(13)} este foarte
mulțumită de ea; este o \VUgras{elevă} bună \textsuperscript{(14)}. Nu ca
haimanaua asta, băiatul meu; numai jocuri are în cap. \VUgras{Învățătorul}
\textsuperscript{(15)} nu îl poate face nicicum să învețe măcar ceva».

%%K t07-tit-4 sub
\VUsaut{3.0mm}
\VUpk{10.2pt}{40}{Vorbele satului.}
\VUsaut{3.0mm}

%%K t07-c1-10-1 p
10. --- Apoi bătrânul ne-a povestit noutățile satului. Aveau de gând să
construiască o școală nouă, fiindcă aceea din \VUgras{primărie} a ajuns prea
strâmtă. Nu ar fi fost greu, căci ajungea să cumpere o parte din grădina
\VUgras{morarului}. Bietului om i-ar fi prins tare bine. De la gerurile cele mari
\VUgras{pietrele} \VUgras{morii} \textsuperscript{(16)} nu au putut măcina grâul,
fiindcă \VUgras{roata} \textsuperscript{(17)} a fost prinsă de \VUgras{gheață}
\textsuperscript{(25)} pe \VUgras{pârâu} \textsuperscript{(23)}.

%%K t07-c1-11-1 p
11. --- «Fiindcă veni vorba de clădiri: ați văzut noua noastră \VUgras{biserică}
\textsuperscript{(18)}? Este foarte pitorească sub căciula de zăpadă.
\VUgras{Ciorile} \textsuperscript{(19)}, care nu își mai recunosc vechea
clopotniță, se așază posomorâte pe copacul cel mare din \VUgras{cimitir}
\textsuperscript{(20)}».

%%K t07-c1-12-1 p
12. --- Pomenirea cimitirului i-a adus aminte cizmarului de cunoscuții tatei care
au murit în anul trecut. «Câte \VUgras{morminte} \textsuperscript{(21)}, domnule,
s-au adăugat celor vechi sub \VUgras{chiparos} \textsuperscript{(22)}! Moartea l-a
luat pe bătrânul nostru notar. Tot satul a fost la \VUgras{înmormântarea} lui».

%%K t07-c1-13-1 p
13. --- «Uite-l pe urmașul lui patinând pe pârâu. A trecut deunăzi ca să îi dreg
cureaua \VUgras{patinei} \textsuperscript{(26)}, care s-a rupt».

%%K t07-c1-14-1 p
14. --- «Iar colo, pe malul pârâului, noul \VUgras{paznic de câmp}
\textsuperscript{(27)}. Este fost jandarm și spaima haimanalelor din sat. Se
împrăștie de îndată ce îi zăresc \VUgras{jambierele} \textsuperscript{(29)} și
\VUgras{șapca} \textsuperscript{(28)}».

%%K t07-c1-15-1 p
15. --- «A! uite-l pe bătrânul Iosif ducând o \VUgras{căruță de vreascuri}
\textsuperscript{(30)} brutarului. --- Într-adevăr, a spus tata, îi recunosc
hamurile de \VUgras{catâri} \textsuperscript{(31)}. Am de vorbit cu el, și mă voi
folosi de prilej. Rămâi cu bine, bătrâne Iacob».

%%K t07-c1-16-1 p
16. --- Tocmai când ieșeam, copiii satului ieșeau din școală și se aruncau pe
\VUgras{ghețușul} \textsuperscript{(33)}, primejduindu-se să cadă și să își rupă
ceva la \VUgras{cădere} \textsuperscript{(34)}. Unul dintre ei și-a luat de la
rotar \VUgras{sania} \textsuperscript{(32)}, a pus în ea un tovarăș și a pornit în
fugă.

%%K t07-c1-17-1 p
17. --- Alții s-au repezit la \VUgras{nămetele} \textsuperscript{(37)} de lângă
\VUgras{pod} \textsuperscript{(40)} și au început să bombardeze \VUgras{omul de
zăpadă} \textsuperscript{(39)}, pe care îl făcuseră în ajun și căruia îi dăduseră
o pălărie, o pipă și o traistă de școală pe umăr.

%%K t07-c1-18-1 p
18. --- Nu le păsa de vântul înghețat și de ger. Cei mai mulți nu aveau nici măcar
\VUgras{fular} \textsuperscript{(35)}, iar ca să se încălzească după jocul cu
\VUgras{bulgări de zăpadă} \textsuperscript{(38)}, le ajungea un pumn de
\VUgras{castane} \textsuperscript{(42)} fierbinți, cumpărate de la bătrâna
Jacqueline, \VUgras{precupeața}, care tremură de frig sub \VUgras{șalul} ei prea
subțire \textsuperscript{(43)}.

%%K t07-tit-5 sub
\VUsaut{3.0mm}
\VUcentre{12.0pt}{0}{\textit{A doua scenă.}}
\VUsaut{3.0mm}

%%K t07-tit-6 sub
\VUpk{10.2pt}{40}{Ferma primăvara.}
\VUsaut{4.0mm}

%%K t07-c2-01-1 p
1. --- Cu toate bucuriile iernii, primim cu veselie adâncă întoarcerea
\VUgras{primăverii}.

%%K t07-c2-01-2 p
Ce tablou mângâietor este curtea fermei seara, în mai!

%%K t07-c2-02-1 p
2. --- La zare \VUgras{soarele} \textsuperscript{(1)} apune încet, scăldând
\VUgras{câmpurile} \textsuperscript{(3)} în \VUgras{raze} purpurii
\textsuperscript{(2)}. \VUgras{Plugarul} harnic \textsuperscript{(6)} se grăbește
să își are \VUgras{ogorul} \textsuperscript{(4)}.

%%K t07-c2-02-2 p
Cu \VUgras{plugul} său \textsuperscript{(5)}, tras de un \VUgras{cal} puternic
\textsuperscript{(7)}, trage \VUgras{brazda} \textsuperscript{(8)}, în care
\VUgras{semănătorul} \textsuperscript{(9)} va împrăștia cu pumnul \VUgras{sămânța}
care va da spice de aur.

%%K t07-c2-03-1 p
3. --- \VUgras{Cosașii} \textsuperscript{(11)} au cosit cu coasele lor iarba
livezii, întorcătorii au uscat \VUgras{fânul} \textsuperscript{(10)},
întorcându-l cu \VUgras{furcile} \textsuperscript{(12)} și cu greblele, și iată-i
că se întorc acasă.

%%K t07-c2-03-2 p
Unii merg înaintea carului greu tras de boi, ducând uneltele pe umăr. Alții, mai
leneși, s-au așezat în tihnă pe fânul înmiresmat.

%%K t07-c2-04-1 p
4. --- Trecând, îl salută prietenește pe \VUgras{grădinar}
\textsuperscript{(16)}, care trudește în \VUgras{livadă} \textsuperscript{(13)},
retezând \VUgras{pomii} și altoind \VUgras{perii} \textsuperscript{(14)}.
\VUgras{Prunii} \textsuperscript{(15)} sunt în floare, iar în \VUgras{grădina de
zarzavat} bine gunoită \textsuperscript{(17)} legumele, varza și salata sunt deja
frumoase.

%%K t07-c2-04-2 p
Pe gardul scund un \VUgras{păun} frumos \textsuperscript{(18)} își desface coada,
ca să își arate penele minunate.

%%K t07-tit-7 sub
\VUpk{10.2pt}{40}{Curtea fermei.}
\VUsaut{3.0mm}

%%K t07-c2-05-1 p
5. --- Ușile \VUgras{grajdului} \textsuperscript{(19)} sunt deschise, și se văd
ieslea și \VUgras{așternutul} \textsuperscript{(23)} \VUgras{iepei}
\textsuperscript{(21)}, pe care \VUgras{îngrijitorul de cai}
\textsuperscript{(20)} tocmai a deshămat-o. \VUgras{Mânzul} \textsuperscript{(22)},
atât de caraghios pe picioarele lui lungi, merge după mamă, zburdând.

%%K t07-c2-06-1 p
6. --- Lângă \VUgras{fântână} \textsuperscript{(29)} \VUgras{vacile}
\textsuperscript{(25)} vin să bea dintr-o \VUgras{albie} de lemn
\textsuperscript{(26)}, care le slujește de adăpătoare. \VUgras{Vițelul}
\textsuperscript{(24)} se folosește de prilej ca să sugă, iar \VUgras{boul}
\textsuperscript{(27)}, simțind mirosul bun al nutrețului, mugește vesel și își
înalță mândru capul cu \VUgras{coarnele} puternice \textsuperscript{(28)}.

%%K t07-c2-07-1 p
7. --- Toți sunt la lucru. \VUgras{Lucrătoarea} \textsuperscript{(32)} ține
\VUgras{funia} \textsuperscript{(31)} fântânii. Scoate apă cu \VUgras{găleata}
\textsuperscript{(30)}, ca să adape vitele. Lângă \VUgras{șură}
\textsuperscript{(33)} un om adună paiele risipite, iar înaintea ușii
\VUgras{fermei} \textsuperscript{(34)} o \VUgras{torcătoare} bătrână
\textsuperscript{(35)} cu furca ei și cu \VUgras{caierul} toarce \VUgras{inul} sau
\VUgras{cânepa}, ca în zilele de altădată.

%%K t07-tit-8 sub
\VUsaut{3.0mm}
\VUpk{10.2pt}{40}{Fermierul.}
\VUsaut{3.0mm}

%%K t07-c2-08-1 p
8. --- \VUgras{Fermierul} \textsuperscript{(36)} cârmuiește toată această muncă și
le dă oamenilor săi îndrumări. Poruncește să fie strânse sub șopron
\VUgras{grapa} \textsuperscript{(40)} și \VUgras{târnăcopul}
\textsuperscript{(39)}, lăsate lângă \VUgras{cușca} \textsuperscript{(38)}
\VUgras{câinelui} \textsuperscript{(37)}.

%%K t07-c2-09-1 p
9. --- Strigătele lui sperie \VUgras{porumbelul} \textsuperscript{(41)}, care
zboară cât poate spre \VUgras{porumbar} \textsuperscript{(42)}, și
\VUgras{găinile} \textsuperscript{(43)}, care se grăbesc înapoi spre
\VUgras{coteț} \textsuperscript{(44)}.

%%K t07-c2-10-1 p
10. --- Înaintea \VUgras{staulului} \textsuperscript{(48)} bătrânul
\VUgras{păstor} \textsuperscript{(45)} își adună \VUgras{turma}
\textsuperscript{(49)}. Ținându-și \VUgras{bâta} \textsuperscript{(46)}, care nu se
desparte de el, ca nici \VUgras{bluza} decolorată \textsuperscript{(47)}, duce în
țarc \VUgras{berbecii} \textsuperscript{(50)}, \VUgras{oile}
\textsuperscript{(53)}, \VUgras{mieii} \textsuperscript{(54)},
\VUgras{caprele} \textsuperscript{(51)} și \VUgras{iezii} \textsuperscript{(52)}.
\VUgras{Câinele} său credincios \textsuperscript{(55)}, Argus, merge cel din urmă
și îi grăbește cu lătratul pe cei rămași în urmă.

%%K t07-c2-11-1 p
11. --- Tot acest zgomot și această forfotă nu tulbură liniștea bunei \VUgras{vaci
cu lapte} \textsuperscript{(56)}, care își zdrăngăne \VUgras{talanga}
\textsuperscript{(57)} în vreme ce este mulsă înaintea ușii \VUgras{grajdului de
vite} \textsuperscript{(58)}.

%%K t07-c2-12-1 p
12. --- Parcă înțelege că trebuie să dea lapte, ca \VUgras{fermiera}
\textsuperscript{(60)} să poată bate \VUgras{untul} în \VUgras{putinei}
\textsuperscript{(61)} sau să facă \VUgras{brânzeturi} alese
\textsuperscript{(64)}, care se zvântă pe o scândură pusă de-a curmezișul
\VUgras{cazanului} \textsuperscript{(63)}.

%%K t07-c2-13-1 p
13. --- Toată \VUgras{lăptăria} \textsuperscript{(59)} o ține curată
\VUgras{argatul} \textsuperscript{(65)}, care tocmai a spălat \VUgras{bidoanele}
\textsuperscript{(62)} în găleata mare pe care o duce.

%%K t07-tit-9 sub
\VUsaut{3.0mm}
\VUpk{10.2pt}{40}{Curtea păsărilor.}
\VUsaut{3.0mm}

%%K t07-c2-14-1 p
14. --- Cu saboții lui grei de lemn umblă destul de stângaci și de aceea sperie
\VUgras{curcanul} \textsuperscript{(68)}, \VUgras{rațele}
\textsuperscript{(66)}, \VUgras{rățoii} \textsuperscript{(67)} și
\VUgras{gâștele} \textsuperscript{(69)}, care își cască larg \VUgras{ciocurile}
\textsuperscript{(70)} și ridică o gălăgie neliniștită.

%%K t07-c2-14-2 p
\VUgras{Cocoșul} \textsuperscript{(71)}, mai războinic, își înalță \VUgras{creasta}
roșie \textsuperscript{(72)}, își scutură penele \VUgras{cozii}
\textsuperscript{(73)} și scoate un «cucurigu» răsunător.

%%K t07-c2-15-1 p
15. --- \VUgras{Cloșca} \textsuperscript{(74)} scurmă în \VUgras{grămada de
gunoi} \textsuperscript{(81)} cu \VUgras{puii} ei \textsuperscript{(75)}.
\VUgras{Purcelul} \textsuperscript{(78)} aleargă la mamă, la \VUgras{scroafa}
uriașă \textsuperscript{(77)}, care se vede lângă coteț.

%%K t07-c2-16-1 p
16. --- În cuștile lor \VUgras{iepurii} \textsuperscript{(79)} mănâncă lacom
\VUgras{iarba} fragedă \textsuperscript{(80)}, pe care le-o aduce fata fermierului.
Eugenia își iubește mult iepurii și în fiecare zi, după ce ia din coteț
\VUgras{ouăle} proaspete \textsuperscript{(76)}, vine să își cerceteze micii
prieteni.
\VUsaut{6.0mm}
\VUfilet{22mm}
